\input{preamble.tex}

\title{\huge{Oblig 2}}
\author{\LARGE{Thobias Høivik}}
\date{}

\begin{document}
\maketitle

\newpage
\begin{prob*}[1]
    We will prove the following statement:

    \[ 
        \text{Every separable Hilbert space} \ H \ \text{has an 
        orthonormal basis.}
    \]

    \begin{enumerate}[label=\alph*)]
        \item Let $D = \{u_1,u_2,\ldots\}$ be a dense subset 
            of $H$ (why does such a set exist?). 
            Show that there is a subset $D' \subseteq D$ with 
            elements $D' = \{v_1,v_2,\ldots\}$, where 
            $(v_1,v_2,\ldots)$ is linearly independent and 
            $\operatorname{span}(v_1,v_2,\ldots)$ is dense in 
            $H$.
        \item Apply the Gram-Schmidt orthogonalization procedure 
            to the list $(v_1,v_2,\ldots)$. Explain why the resulting 
            list $(e_1,e_2,\ldots)$ is orthonormal and its span 
            is dense in $H$.
        \item Prove that $(e_1,e_2,\ldots)$ is an orthonormal 
            basis for $H$.
    \end{enumerate}
\end{prob*}
\begin{proof}[Proof of (a)]
    Trivially, every hilbert space $H$ is dense in itself since 
    $\overline H = H$. In our case we are given the assumption 
    that $H$ is separable, so it follows that there exists 
    a countable dense $D \subseteq H$.

    We will construct $D' \subseteq D$ with the desired properties. 
    Let $v_1$ be the first element of $D$ which is non-zero. 
    Suppose we have chosen $n$ elements $D_n' = \{v_1,\ldots,v_n\}$. 
    Let $D_{n+1}' = D_n' \cup \{u_i \in D\setminus D_n'\}$ where 
    $u_i$ is the first element of $D \setminus D_n'$ which is 
    not in $\operatorname{span}(D_n')$. Let 
    \[
        D' := \bigcup_{n \in \mathbb N} D'_{n}.
    \]
    Then $D'$ is linearly independent by construction.

    Now $D \subseteq \operatorname{span}(D')$ because if 
    $u_i \not\in D'$ then 
    $u_i \in \operatorname{span}(D'_n)$ for some $n$, so 
    $u_i \in \operatorname{span}(D')$. If $u_i \in D'$ the inclusion 
    is trivial. 
    Now we see that 
    \[
        H = \overline D \subseteq \overline{\operatorname{span}(D')} 
        \subseteq H,
    \]
    so $D'$ is dense in $H$.
\end{proof}

\begin{proof}[Proof of (b)]
    Apply the Gram-Schmidt procedure to the linearly independent list
    $(v_1,v_2,\ldots)$. Define
    \[
        w_1 = v_1, \qquad e_1 = \frac{w_1}{\|w_1\|},
    \]
    and recursively
    \[
        w_n = v_n - \sum_{k=1}^{n-1} \langle v_n,e_k\rangle e_k,
        \qquad
        e_n = \frac{w_n}{\|w_n\|}.
    \]
    Since $(v_1,v_2,\ldots)$ is linearly independent, we have
    $w_n \neq 0$ for every $n$, so each $e_n$ is well-defined.

    By the Gram-Schmidt construction, the vectors $e_1,e_2,\ldots$
    are orthonormal. Moreover, for every $n$,
    \[
        \operatorname{span}(e_1,\ldots,e_n)
        =
        \operatorname{span}(v_1,\ldots,v_n).
    \]
    Hence
    \[
        \operatorname{span}(e_1,e_2,\ldots)
        =
        \operatorname{span}(v_1,v_2,\ldots).
    \]
    But by part (a), $\operatorname{span}(v_1,v_2,\ldots)$ is dense
    in $H$. Therefore $\operatorname{span}(e_1,e_2,\ldots)$ is also
    dense in $H$.
\end{proof}

\begin{proof}[Proof of (c)]
    We proved orthonormality in (b), and furthermore we showed 
    $H = \operatorname{span}(D')$ in (a). 
    Since $\operatorname{span}(D') = 
    \operatorname{span}(e_1,e_2\ldots)$
    we conclude that $(e_1,e_2,\ldots)$.
\end{proof}

\begin{prob*}[2]
    We will prove the Riesz Representation Theorem:

    \begin{quote}
        If \(H\) is a Hilbert space over \(\mathbb{R}\) and
        \(T \in \mathcal{L}(H,\mathbb{R})\), then there exists a unique
        vector \(v \in H\) such that
        \[
            Tu = \langle u,v\rangle
            \quad \text{for all } u \in H.
        \]
    \end{quote}

    We assume $H$ is separable, so an orthonormal basis 
    $(e_k)_{k\in\mathbb N}$.

    \begin{enumerate}[label=\alph*)]
        \item Show that there is some $C > 0$ such that 
            \[
                |Tu| \leq C \left(\sum_{k=1}^\infty |\langle 
                u, e_k\rangle|^2\right)^{1/2} \ 
                \forall u \in H.
            \]
        \item Define $\alpha_k := Te_k$ for $k \in \mathbb N$. 
            Show that 
            \[
                \sum_{k=1}^n|\alpha_k|^2 \leq 
                C\left(\sum_{k=1}^n|\alpha_k|^2\right)^{1/2} 
                \ 
                \forall n \in \mathbb N.
            \]
            Conclude that $\alpha = (\alpha_1,\alpha_2,\ldots) \in 
            \ell^2(\mathbb R)$ and that $\|\alpha\|_{\ell^2} \le C$.
        \item Show that the series $\sum_{k=1}^\infty \alpha_k e_k$
            converges, and that 
            \[
                Tu = \langle u,v\rangle \ \forall u \in H.
            \]
    \end{enumerate}
\end{prob*}
\begin{proof}[Proof of (a)]
    $T$ is bounded (so there 
    exists $C > 0$ so $|Tu| \leq C\|u\|$) and combining 
    this with Parseval's identity the desired result is immediate.
\end{proof}
\begin{proof}[Proof of (b)]
    Let $\beta_n := \sum_{k=1}^n \alpha_ke_k$. Then 
    \[
        T\beta_n = T \left(\sum_{k=1}^n\alpha_ke_k\right) 
        = \sum_{k=1}^n\alpha_k T e_k = \sum_{k=1}^n \alpha_k^2.
    \]
    Since $\alpha_k$ is real we have $|\alpha_k|^2 = \alpha_k^2$.
    Hence 
    \[
        T\beta_n = \sum_{k=1}^n |\alpha_k|^2.
    \]
    Applying (a) we get 
    \[
        \sum_{k=1}^n |\alpha_k|^2 = 
        |T\beta_n| \le C \left( 
        \sum_{j=1}^n|\langle \beta_n, e_j \rangle|^2\right)^{1/2}.
    \]
    Since $\beta_{j} \ (j > n)$ is $0$, while 
    $\langle \beta_j, e_j \rangle = a_j$ otherwise we get the 
    sum equaling $C \sum_{k=1}^n |\alpha_k|^2$. So 
    \[
        \sum_{k=1}^n |\alpha_k|^2 \leq 
        C \left(\sum_{k=1}^n |\alpha_k|^2\right)^{1/2}.
    \]
    Now doing some rearranging we can divide both sides by 
    the square root of the sum to get 
    \[
        \sqrt{\sum_{k=1}^n |\alpha_k|^2} \leq C.
    \]
    Then $\sum_{k=1}^n(|\alpha_k|^2)^{1/2} \leq C^2$ so all finite 
    sums are bouded, increasing and nonnegative, hence 
    \[
        \|\alpha\|_{\ell^2} \leq C < \infty.
    \]
\end{proof}

\begin{prob*}[3]
    Let $X$ be a set and let $S := \{\{x\} : x \in X\}$ (the 
    set of singletons in $X$). Let $\mathcal A := \sigma(S)$ be 
    the $\sigma$-algebra generated by $S$.

    \begin{enumerate}[label=\alph*)]
        \item Show that 
            \[
                \mathcal A = \{E \subseteq X : E \ 
                \text{or} \ E^c \ \text{is countable}\}.
            \]
        \item Assume that $X$ is uncountable. Show that there exists
            a set $E \subseteq X$ which is not meassurable with 
            respect to $\mathcal A$.
    \end{enumerate}
\end{prob*}
\begin{proof}[Proof of (a)]
    Define $F := \{E \subseteq X : E \ \text{or} \ E^c \ 
    \text{is countable}\}$.
    
    $S \subseteq F$ is trival.
    Let $E \in F$. If $E$ is countable, then $E^c$ is co-countable, 
    so $E^c \in F$. Likewise if we suppose $E^c$ is countable.
    Let $E_1,E_2,\ldots \in F$. If all $E_n$ are countable then 
    the countable union is countable (the proof of this 
    is straightforward), so $\cup_{i=1}^\infty E_i \in F$. Otherwise 
    some $E_k$ is co-countable. Then 
    $(\cup_{i=1}^\infty E_i) = \cap_{i=1}^\infty E_i^c \subseteq 
    E_k^c$. Since $E_k^c$ is countable we have 
    $\cup_{n=1}^\infty E_i \in F$.
    So $F$ is a $\sigma$-algebra containing $S$. Hence by 
    minimality: 
    \[
        \sigma(S) = \mathcal A \subseteq F.
    \]

    Every singleton is in $S \subseteq \mathcal A$. Any 
    countable $E = \{x_1,x_2,\ldots\}$ can be expressed as 
    a countably infinite union of singletons so $E \in \mathcal A$.
    If $E^c$ is countalbe, then $E = (E^c)^c \in \mathcal A$.
    So all of $F$ lies in $\mathcal A$.
    
    Both inclusions yield 
    \[
        \mathcal A = F.
    \]
    
\end{proof}
\begin{proof}[Proof of (b)]
    If we for instance take 
    $X = \mathbb R$ and consider $E = (0,\infty)$. Then 
    $E$ is uncountable and $E^c = (-\infty,0]$ is also uncountable. 
    Then $E, E^c \not\in \mathcal A$.

    For arbitrary $X$ with $|X| \geq \aleph_1$ we need 
    the fact that such a set we can necessarily pick out an 
    uncountable set such that the complement is still uncountable. 
\end{proof}

\begin{prob*}[4]
    Let $X = \{x_1,\ldots,x_n\}$ be a set containing 
    $n \in \mathbb N$ elements, and equip it with the discrete 
    $\sigma$-algebra $\mathcal A = 2^X$. Let $p_1,\ldots,p_n \in [0,1]$
    be such that $\sum_{k=1}^n p_k = 1$. Show that there exists 
    a unique probability meassure $\mu$ on $(X,\mathcal A)$ with 
    the property that $\mu(\{x_k\}) = p_k$ for $k = 1,\ldots,n$.
\end{prob*}
\begin{proof}
    It is clear that if we define, for $E \subseteq X$ (equivalently
    $E \in 2^X$), 
    \[
        \mu(E) := \sum_{x_k \in E} p_k.
    \]
    
    Suppose $\nu$ are two probability meassures with 
    $\mu(\{x_k\}) = p_k = \nu(\{x_k\})$. 
    Take any $E \in 2^X$. Since $X = \{x_1,\ldots,x_n\}$, we can 
    write 
    \[
        E = \bigcup_{x_k \in E} \{x_k\}, 
    \]
    as a finite disjoint union. By additivity 
    \[
        \mu(E) = \sum_{x_k \in E}\mu(\{x_k\}) = \sum_{x_k \in E}p_k,
    \]
    and similarly for $\nu$, 
    \[
        \nu(E) = \cdots = \sum_{x_k \in E}p_k. 
    \]
    We therefore conclude that 
    \[
        \mu(E) = \nu(E).
    \]
    As $E$ was chosen arbitrarily we conclude that $\mu$ and 
    $\nu$ agree on all meassureable sets.
\end{proof}

\begin{prob*}[5]
    We toss a coin three times, and record the result as a 
    vector $(a_1,a_2,a_3)$, where $a_1,a_2,a_3 \in \{0,1\}$, where 
    $0$ encodes "heads" and $1$ encodes "tails". Let 
    \[
        X := \{(a_1,a_2,a_3) : a_1,a_2,a_3 \in \{0,1\}\},
    \]
    and let $\mathcal A := 2^X$.

    \begin{enumerate}[label=\alph*)]
        \item If heads and tails have equal likelyhoods, what is the 
            correct probability meassure on the meassureable 
            space $(X,\mathcal A)$?
        \item Find the set $E \in \mathcal A$ which describes the 
            event "at least two tosses were tails".
        \item Two events $E,F \in \mathcal A$ are independent if 
            the probability of both occuring equals the 
            product of the probability of each of them: 
            \[
                \mu(E \cap F) = \mu(E)\mu(F).
            \]
            Show that the events "the first toss is heads" and 
            "the first and second tosses are the same" are 
            independent.
    \end{enumerate}
\end{prob*}
\begin{proof}[Solution of (a)]
    If heads and tails have equal likelyhoods then 
    the chance of any sequence is the same. Since 
    $|X| = 2^3 = 8$ we then want 
    \[
        \mu(\{\omega\}) = \frac{1}{8} 
    \]
    for any outcome $\omega \in X$.
    We define, for $E \in \mathcal A$, 
    \[
        \mu(E) = \frac{|E|}{8}.
    \]
\end{proof}
\begin{proof}[Solution of (b)]
    \[
        E = \{(1,1,1),(1,1,0),(1,0,1),(0,1,1)\}.
    \]
\end{proof}
\begin{proof}[Proof of (c)]
    Let $E$ be the set of events where the first toss is heads 
    and let $F$ be the set of events where the first and second 
    tosses are the same. Then 
    \[
        E = \{(0,a_2,a_3) : a_2,a_3 \in \{0,1\}\},
    \]
    and 
    \[
        F = \{(a_1,a_2,a_3) : a_1 = a_2, \ a_1,a_2,a_3 \in \{0,1\}\}.
    \]
    In particular, 
    \[
        E = \{(0,0,0), (0,0,1), (0,1,0), (0,1,1)\}, 
    \]
    and 
    \[
        F = \{(1,1,1), (0,0,0), (1,1,0), (0,0,1)\}.
    \]
    Then the meassure of both is $1/2$. Their intersection 
    is size $2$ so the meassure 
    \[
        \mu(E\cap F) = 2/8 = 1/4,
    \]
    and 
    \[
        \mu(E)\mu(F) = 1/2 \cdot 1/2 = 1/4.
    \]
    Thus we conclude that these events are independent.
\end{proof}

\begin{prob*}[6]
    Let $(X,\mathcal A,\mu)$ be a meassure space. For a list of 
    meassurable sets $E_1,E_2,\ldots \in \mathcal A$, we define 
    the limit superior 
    \[
        \underset{n\to\infty}{\lim\sup}E_n := 
        \bigcap_{n=1}^\infty \bigcup_{k=n}^\infty E_k.
    \]

    \begin{enumerate}[label=\alph*)]
        \item Prove that $\lim\sup_{n\to\infty}E_n$ is meassurable.
        \item Prove that 
            \begin{align*}
                \underset{n\to\infty}{\lim\sup}E_n 
                \\ = \{x \in X : x \ \text{lies in infinitely 
                many of the sets} \ E_1,E_2,\ldots\}.
            \end{align*}
        \item Prove the Borel-Cantelli lemma: 
            \[
                \text{If} \ \sum_{k=1}^\infty \mu(E_k) < \infty, 
                \ \text{then} \ 
                \mu(\lim\sup_{n\to\infty} E_n) = 0.
            \]
    \end{enumerate}
\end{prob*}
\begin{proof}[Proof of (a)]
    We have a list $E_1,E_2,\ldots \in \mathcal A$.

    We notice that 
    \[
        \bigcap_{n=1}^\infty 
        \underbrace{\bigcup_{k=n}^\infty E_k}_{\text{meassurable}}, 
    \]
    since that part is a countably infinite union of meassurable 
    sets. Taking the countable union of meassurable sets 
    is meassurable, hence 
    \[
        \underset{n\to\infty}{\lim\sup}E_n
    \]
    is meassurable.
\end{proof}
\begin{proof}[Proof of (b)]
    Let $x \in \underset{n\to\infty}{\lim\sup}E_n$.
    Then $\forall n \in \mathbb N$, $x \in \cup_{k=n}^\infty E_k$, 
    meaning $\forall n \in \mathbb N, \ \exists k \geq n$ such 
    that $x \in E_k$. Therefore $x$ lies in infinitely many 
    $E_k$. 

    Conversely, if $x$ lies in infinitely many $E_k$, then for 
    every $n$ there is some $k \geq n$ with $x \in E_k$, so 
    $x \in \cup_{k=n}^\infty E_k$ for every $n$. 
    Hence $x \in \underset{n\to\infty}{\lim\sup}E_n$.
\end{proof}
\begin{proof}[Proof of (c)]
    $\mu(E_k) \geq 0$ so since the sum of all $\mu(E_k)$ is finite 
    we conclude that 
    \[
        \sum_{k=n}^\infty \mu(E_k) \to 0.
    \]
    Then 
    \[
        \mu(\underset{n\to\infty}{\lim\sup}E_n) 
        \leq 
        \mu \left( \bigcup_{k=n}^\infty E_k \right)
        \leq 
        \sum_{k=n}^\infty \mu(E_k) \to 0.
    \]
\end{proof}


\end{document}
